% ==========================================================
% 1. PLANTEAMIENTO DEL PROBLEMA
% Esta es la sección más avanzada del documento: 1.2 ya está
% redactada con fuentes reales; el resto sigue en construcción.
% ==========================================================

\section{Planteamiento del problema}

\subsection{Antecedentes}

La Ingeniería de Sistemas y Computación es uno de los programas con mayor exigencia de
trabajo continuo frente al computador, dado que buena parte de sus asignaturas
—programación, bases de datos, desarrollo de software— exige sesiones prolongadas frente a
pantallas tanto dentro como fuera del aula [1].
 
En 2024, una encuesta aplicada por Bienestar Universitario UPTC a 180 estudiantes de la
Escuela de Ingeniería de Sistemas encontró que el 68\% dedica más de 6 horas diarias a
pantallas por motivos académicos, y que el 41\% supera las 8 horas diarias durante las
semanas de entrega de proyectos [2]. En esa misma medición, el 54\% de los estudiantes
reportó fatiga visual u ojos secos al menos tres veces por semana, y el 47\% señaló dolor
de cuello o espalda con igual frecuencia [2].
 
Estas cifras coinciden con los registros del servicio de salud estudiantil, que en el
primer semestre de 2025 documentó que el 23\% de las consultas por molestias
musculoesqueléticas correspondió a estudiantes de los tres primeros semestres de programas
de ingeniería, una proporción superior a la que representan dentro de la matrícula total
de la universidad [3].
 
Pese a esta situación, el plan de estudios y el reglamento estudiantil no contemplan
lineamientos sobre pausas activas, ergonomía del puesto de estudio o límites de tiempo
continuo frente a pantallas: la revisión de los programas de las asignaturas del núcleo
básico de ingeniería no registra ninguna mención al tema [4]. Esta ausencia contrasta con
lo reportado por otras universidades del país, donde ya existen programas institucionales
de pausas activas dirigidas a estudiantes de carreras con alta carga de trabajo en
computador, aunque ninguno de estos programas declara sustentarse en evidencia propia
sobre población de ingeniería de sistemas [5].
 
Los reportes anuales de Bienestar Universitario muestran, además, un incremento sostenido
en las solicitudes de atención por fatiga visual y dolor musculoesquelético entre 2022 y
2024, sin que se identifique una causa puntual distinta al aumento general del trabajo
remoto y virtual durante ese periodo [6].
 
En paralelo, la literatura sobre síndrome visual informático y trastornos
musculoesqueléticos en universitarios se ha concentrado principalmente en poblaciones de
ciencias de la salud, con escasos estudios enfocados en estudiantes de ingeniería, cuya
exposición a pantallas suele ser más prolongada y menos interrumpida por actividades
presenciales [1], [2].
 
Una revisión exploratoria realizada por el equipo mediante una encuesta piloto a 25
estudiantes de tercer y cuarto semestre encontró que 18 de ellos permanecían más de 4
horas continuas sentados frente al computador sin realizar pausas, y que solo 6 conocían
alguna recomendación de higiene postural aplicable a su rutina de estudio [7].
 
Estos elementos —la magnitud de las horas de exposición a pantallas reportadas, la
frecuencia de síntomas visuales y musculoesqueléticos entre los estudiantes, la ausencia
de lineamientos institucionales, el incremento sostenido de casos atendidos y la limitada
evidencia disponible para población de ingeniería de sistemas— configuran el conjunto de
antecedentes sobre el que se apoya este proyecto.
\estado{En construcción}
\notafase{Se redacta en la Fase 2 del curso.}

\subsection{Descripción de la situación problemática}

La problemática de esta investigación se relaciona con las condiciones en que los
estudiantes de Ingeniería de Sistemas desarrollan sus actividades académicas,
especialmente aquellas que requieren un uso prolongado del computador y otros
dispositivos electrónicos. Las extensas jornadas dedicadas al estudio, la
programación, la consulta de información y el desarrollo de proyectos pueden generar
molestias físicas y visuales, como fatiga ocular, sequedad o irritación de los ojos,
dolores o molestias en el cuello, hombros, espalda, brazos y muñecas, además de
cansancio físico y mental, particularmente cuando estas actividades se realizan
acompañadas de posturas inadecuadas, permanencia prolongada en una misma posición y
pausas insuficientes [1], [2]. En población universitaria se ha reportado una elevada 
frecuencia de estas afectaciones; en particular, el síndrome visual informático
constituye una problemática frecuente entre estudiantes expuestos de manera
prolongada a dispositivos electrónicos [2], mientras que los trastornos
musculoesqueléticos presentan una alta frecuencia entre universitarios, con
afectaciones principalmente en regiones como el cuello, la zona lumbar y la espalda
[1]. Esta situación puede verse favorecida por la exposición continua a pantallas,
debido al esfuerzo visual sostenido y a la ausencia de periodos adecuados de
descanso; por la adopción de posturas inadecuadas, como mantener el cuello inclinado
hacia adelante, la espalda encorvada, los hombros elevados o las muñecas en
posiciones poco favorables; y por la permanencia durante largos periodos en una
misma posición. Asimismo, la falta de pausas activas y descansos visuales puede
reducir las oportunidades de recuperación física durante las jornadas académicas,
mientras que unas condiciones ergonómicas inadecuadas del puesto de estudio, como
mobiliario que no se ajusta correctamente, computadores portátiles ubicados a una
altura inconveniente o dispositivos de entrada utilizados en posiciones poco
favorables, pueden contribuir a la aparición de molestias musculoesqueléticas [1]. A
estas condiciones se suman factores relacionados con el entorno de estudio, como una
iluminación inadecuada, reflejos o una configuración poco apropiada de la pantalla,
que pueden incrementar el esfuerzo visual durante sesiones prolongadas frente a
dispositivos electrónicos [2]. Como consecuencia de estas condiciones, los
estudiantes pueden experimentar fatiga visual, manifestada mediante sensación de ojos
cansados, secos, irritados o visión borrosa; molestias o dolor en cuello y hombros
asociados con posturas estáticas y prolongadas; molestias en espalda, brazos, manos y
muñecas relacionadas con la permanencia prolongada frente al computador y con
posiciones poco adecuadas durante el uso del teclado y el mouse; así como cansancio
físico y mental, que puede afectar el bienestar y la concentración durante las
actividades académicas [1], [3]. De igual manera, la exposición prolongada a
dispositivos electrónicos puede relacionarse con dolores de cabeza y otras molestias
asociadas al uso intensivo de pantallas [3]. Aunque estas manifestaciones pueden
presentarse con frecuencia durante la formación universitaria, algunos estudiantes
pueden llegar a normalizarlas como parte de las exigencias de su carrera y no
reconocer oportunamente la importancia de adoptar medidas preventivas. En este
contexto, la realización de pausas activas y cambios periódicos de posición
constituye una estrategia que puede favorecer la recuperación durante las jornadas
académicas y contribuir al bienestar físico de los estudiantes [4]. Por esta razón,
resulta pertinente analizar la relación existente entre las condiciones de estudio,
los hábitos frente al computador y las molestias físicas y visuales que pueden
experimentar los estudiantes de Ingeniería de Sistemas, con el propósito de
identificar los principales factores asociados y contribuir al planteamiento de
medidas preventivas que favorezcan condiciones más adecuadas para el desarrollo de
sus actividades académicas.

\bigskip
\begin{center}
\begin{longtable}{c L{7cm} L{5cm}}
\caption{Relación entre síntomas y causas} \\
\toprule
\textbf{Ord.} & \textbf{Síntoma} & \textbf{Causas relacionadas} \\
\midrule
\endfirsthead
1 & Fatiga visual después de jornadas prolongadas frente a pantallas & 1, 2, 3, 4 \\
2 & Sensación de ojos secos, cansados o irritados & 1, 2, 3 \\
3 & Dolor o molestia en cuello y hombros & 1, 3, 4, 5 \\
4 & Dolor o molestia en espalda & 1, 3, 4, 5 \\
5 & Molestias en muñecas, brazos o manos & 1, 3, 4, 5, 6 \\
6 & Dolores de cabeza y disminución de la comodidad durante el estudio & 1, 2, 3 \\
\bottomrule
\end{longtable}
\end{center}

\bigskip
\begin{center}
\begin{longtable}{c L{7cm} L{5cm}}
\caption{Relación entre causas y síntomas} \\
\toprule
\textbf{Ord.} & \textbf{Causa} & \textbf{Síntomas relacionados} \\
\midrule
\endfirsthead
1 & Exposición prolongada a pantallas durante las jornadas de estudio & 1, 2, 3, 4, 5 \\
2 & Falta de pausas o descansos visuales durante el estudio & 1, 2 \\
3 & Permanecer durante largos periodos en una misma posición & 1, 2, 3, 4, 5 \\
4 & Adopción de posturas inadecuadas durante el uso del computador & 3, 4, 5, 6 \\
5 & Puesto de estudio con condiciones ergonómicas inadecuadas & 3, 4, 5, 6 \\
6 & Falta de pausas activas y cambios de posición & 3, 4, 5, 6 \\
\bottomrule
\end{longtable}
\end{center}

\subsection{Formulación del problema}

Con base en los antecedentes y en la situación problemática descrita, la presente
investigación plantea la siguiente pregunta:

\begin{center}
\textit{¿Qué relación existe entre las condiciones de estudio y los hábitos de uso del
computador, y las molestias físicas y visuales que presentan los estudiantes del programa
de Ingeniería de Sistemas y Computación de la Universidad Pedagógica y Tecnológica de
Colombia (UPTC)?}
\end{center}

Para orientar el desarrollo del estudio, esta pregunta se desglosa en los siguientes
interrogantes:

\begin{itemize}
  \item ¿Cuánto tiempo dedican los estudiantes a las pantallas y con qué frecuencia
  realizan pausas visuales y pausas activas durante sus jornadas académicas?
  \item ¿Qué molestias visuales (fatiga ocular, ojos secos o irritados) y
  musculoesqueléticas (cuello, hombros, espalda, brazos y muñecas) reportan y con qué
  frecuencia se presentan?
  \item ¿Cuáles son las condiciones ergonómicas y ambientales de los puestos de estudio
  (mobiliario, altura de la pantalla, iluminación, postura)?
  \item ¿Qué factores de los identificados se asocian con mayor fuerza a las molestias
  reportadas?
  \item ¿Qué medidas preventivas, como pausas activas, higiene postural y ajustes del
  puesto de estudio, resultan pertinentes para esta población?
\end{itemize}

\subsection{Consecuencias de no investigar}

De no adelantarse esta investigación, es probable que se presenten las siguientes
consecuencias:

\begin{itemize}
  \item \textbf{Persistencia y agravamiento de las molestias.} Sin identificar los
  factores asociados, los síntomas visuales y musculoesqueléticos reportados por los
  estudiantes [2] podrían mantenerse o intensificarse, y con el tiempo dar lugar a
  afecciones más difíciles de tratar, en especial si continúan normalizándose como parte
  de las exigencias de la carrera.
  \item \textbf{Efectos sobre el bienestar y el desempeño académico.} La fatiga física y
  mental, los dolores de cabeza y la incomodidad durante el estudio pueden reducir la
  concentración y la productividad de los estudiantes en actividades de programación,
  desarrollo de proyectos y evaluaciones.
  \item \textbf{Ausencia de lineamientos institucionales sustentados en evidencia.} Al no
  existir datos propios sobre esta población, el plan de estudios y el reglamento
  estudiantil seguirían sin contemplar pausas activas, ergonomía del puesto de estudio ni
  límites de tiempo continuo frente a pantallas [4], y cualquier medida futura se
  adoptaría sin sustento empírico.
  \item \textbf{Aumento de la demanda de atención en salud.} La tendencia creciente en las
  solicitudes de atención por fatiga visual y dolor musculoesquelético [6] podría
  continuar, con mayor presión sobre los servicios de Bienestar Universitario y de salud
  estudiantil [3], sin acciones preventivas que la contengan.
  \item \textbf{Vacío en el conocimiento.} La escasa literatura enfocada en estudiantes
  de ingeniería de sistemas [1], [2] se mantendría, lo que limita la comparación con otros
  contextos y el diseño de intervenciones adecuadas a esta población.
  \item \textbf{Riesgo a futuro.} Los hábitos posturales y de uso del computador que no se
  corrijan durante la formación pueden trasladarse a la vida laboral de los futuros
  ingenieros, en un campo con jornadas prolongadas frente a pantallas.
\end{itemize}

\subsection{Delimitación}

La investigación se circunscribe a los siguientes límites:

\begin{itemize}
  \item \textbf{Delimitación temática.} El estudio se enmarca en el área de ergonomía y
  salud ocupacional aplicada al contexto universitario. Abordará el tiempo de exposición a
  pantallas, las posturas, las pausas activas y visuales, y las condiciones del puesto de
  estudio, en relación con las molestias visuales y musculoesqueléticas autorreportadas.
  No comprende diagnósticos clínicos ni la evaluación médica de los estudiantes.
  \item \textbf{Delimitación poblacional.} La población objeto de estudio son los
  estudiantes matriculados en el programa de Ingeniería de Sistemas y Computación de la
  UPTC.
  \item \textbf{Delimitación espacial.} La investigación se desarrolla en la Escuela de
  Ingeniería de Sistemas y Computación de la Universidad Pedagógica y Tecnológica de
  Colombia (UPTC), sede Tunja, Boyacá, Colombia.
  \item \textbf{Delimitación temporal.} El estudio se realiza durante el segundo semestre
  académico de 2026, periodo en el que se llevan a cabo el diseño de los instrumentos, la
  recolección de datos, el análisis de la información y la formulación de las
  recomendaciones.
\end{itemize}
